\chapter*{Abstract}

This Bachelor's Thesis presents OmniLitter, a multiplatform system aimed at the
management, collection and analysis of data on litter pollution in terrestrial
and aquatic environments. The project addresses the need to digitalise field
sampling workflows and to unify the information collected on site so that
observations can be queried, traced and reused according to scientific criteria.

The proposed solution integrates a mobile application, a web portal and a
centralised API within a common architecture. This structure covers the full work
cycle, from campaign organisation and field data capture to consultation,
analysis and export of the collected information. The system relies on a shared
litter taxonomy, geolocated observations and synchronisation mechanisms that
support work in contexts with limited connectivity.

The work covers requirements analysis, architectural and data model design,
development planning, implementation of the main modules, deployment and
validation through testing. The result is a functional prototype that
demonstrates the feasibility of a unified platform for research, environmental
monitoring and the assessment of litter reduction measures. OmniLitter therefore
provides an extensible technological foundation on which new capabilities can be
incorporated in future work.

\vspace{.5cm}

\textbf{Keywords:} OmniLitter, litter monitoring, mobile application, web
portal, geolocation, offline work, synchronisation, environmental data analysis.
